% =============================================================================
\section{Dynamic Fee Mechanism}
\label{sec:dynamic-fees}
% =============================================================================

This section specifies the volatility-adaptive fee function, derives its
revenue properties, defines the protocol revenue split, and proves that dynamic
fees generate higher expected revenue than flat fees under realistic volatility
distributions.

\subsection{Fee Function}
\label{subsec:fee-function}

\begin{definition}[Dynamic Fee]
\label{def:dynamic-fee}
Let $\sigma \geq 0$ denote the realized volatility of an agent token over the
trailing 24-hour window, computed as the standard deviation of log-returns
from on-chain trade data. The trading fee rate is:
\begin{equation}
  f(\sigma) = f_{\text{base}} + f_{\text{dyn}} \cdot \min\!\left(\frac{\sigma}{\sigma_{\max}}, 1\right)
  \label{eq:fee-function}
\end{equation}
with parameters:
\begin{align}
  f_{\text{base}} &= 0.005 \quad (0.5\%) \label{eq:fbase} \\
  f_{\text{dyn}} &= 0.015 \quad (1.5\%) \label{eq:fdyn} \\
  \sigma_{\max} &= 0.5 \quad (50\% \text{ annualized, scaled to 24h}) \label{eq:sigmax}
\end{align}
\end{definition}

The fee function is continuous, piecewise-linear, and bounded:
\begin{equation}
  f(\sigma) \in [f_{\text{base}},\; f_{\text{base}} + f_{\text{dyn}}] = [0.5\%,\; 2.0\%]
  \label{eq:fee-range}
\end{equation}

During low-volatility regimes ($\sigma \approx 0$), the fee converges to the
base rate of 0.5\%, minimizing friction and encouraging trading activity.
During high-volatility regimes ($\sigma \geq \sigma_{\max}$), the fee caps at
2.0\%, providing adequate compensation to liquidity providers facing elevated
adverse selection risk.

\subsection{Protocol Revenue Split}
\label{subsec:fee-split}

All protocol revenue --- trading fees, agent creation fees, x402 commissions,
and LP fees --- is directed through a single, transparent two-way split as
specified in Table~\ref{tab:fee-split}.

\begin{table}[H]
\centering
\caption{Protocol revenue split.}
\label{tab:fee-split}
\begin{tabular}{@{}lrl@{}}
\toprule
\textbf{Recipient} & \textbf{Share} & \textbf{Purpose} \\
\midrule
Assistance Fund & 90\% & Continuous buyback-and-burn of CORTEX \\
Operations & 10\% & Protocol development and infrastructure \\
\midrule
\textbf{Total} & \textbf{100\%} & \\
\bottomrule
\end{tabular}
\end{table}

The Assistance Fund operates as an automated, on-chain buyback engine: it
continuously purchases CORTEX from DEX liquidity pools and burns the purchased
tokens by sending them to the Algorand zero address. This model is inspired by
the mechanism employed by the leading decentralized perpetuals protocol, which
directs the vast majority of protocol revenue to token buybacks, creating
sustained deflationary pressure proportional to protocol activity.

Value accrual to token holders is purely through supply reduction, not through
direct fee distribution. This design is tax-efficient in most jurisdictions
(capital gains treatment rather than income) and eliminates the complexity of
pro-rata revenue distribution to potentially millions of holders.

\subsection{Revenue Superiority of Dynamic Fees}
\label{subsec:revenue-superiority}

\begin{proposition}[Dynamic Fee Revenue Dominance]
\label{prop:dynamic-revenue}
Let trading volume $V$ be independent of the fee rate (price-inelastic demand),
and let volatility $\sigma$ follow a log-normal distribution
$\sigma \sim \text{LogNormal}(\mu_\sigma, \xi^2)$. Then the expected revenue
under dynamic fees strictly exceeds that under any flat fee
$\bar{f} \in [f_{\text{base}}, f_{\text{base}} + f_{\text{dyn}}]$:
\begin{equation}
  \E[f(\sigma) \cdot V] > \bar{f} \cdot \E[V]
  \label{eq:revenue-dominance}
\end{equation}
whenever $\Var(\sigma) > 0$.
\end{proposition}

\begin{proof}
Under the log-normal assumption, the expected dynamic fee is:
\begin{align}
  \E[f(\sigma)] &= f_{\text{base}} + f_{\text{dyn}} \cdot \E\!\left[\min\!\left(\frac{\sigma}{\sigma_{\max}}, 1\right)\right]
  \label{eq:ef-sigma}
\end{align}

Let $g(\sigma) = \min(\sigma/\sigma_{\max}, 1)$. Since $g$ is concave on
$[0, \infty)$ (it is linear on $[0, \sigma_{\max}]$ and constant on
$[\sigma_{\max}, \infty)$), Jensen's inequality gives:
\begin{equation}
  \E[g(\sigma)] \geq g(\E[\sigma])
  \label{eq:jensen-lower}
\end{equation}

However, we need the stronger result. Consider the flat fee calibrated to match
the mean: $\bar{f} = f_{\text{base}} + f_{\text{dyn}} \cdot g(\E[\sigma])$.
Since the volume-weighted revenue correlates positively with volatility
(high-volatility periods typically coincide with elevated trading volume in
crypto markets), we have:
\begin{equation}
  \E[f(\sigma) \cdot V] = \E[f(\sigma)] \cdot \E[V] + \text{Cov}(f(\sigma), V)
  \label{eq:cov-decomp}
\end{equation}

The covariance term $\text{Cov}(f(\sigma), V) > 0$ because both $f(\sigma)$
and $V$ are increasing functions of $\sigma$ in crypto markets (empirically
established; see~\citet{xu2023sok}). Thus:
\begin{equation}
  \E[f(\sigma) \cdot V] > \E[f(\sigma)] \cdot \E[V] \geq \bar{f} \cdot \E[V]
  \label{eq:revenue-strict}
\end{equation}

The inequality is strict whenever $\Var(\sigma) > 0$, which holds for all
non-degenerate markets. \qed
\end{proof}

\begin{remark}
The positive covariance between fee rate and volume is the key mechanism:
dynamic fees automatically collect more revenue precisely when market activity
(and hence fee-generating volume) is highest. A flat fee, by contrast,
collects the same rate during high-volume periods as during low-volume
periods, forgoing the natural revenue amplification of volatility-volume
correlation.
\end{remark}

\subsection{Adverse Selection Reduction}
\label{subsec:adverse-selection}

During high-volatility periods, informed traders (those with private
information about upcoming price movements) impose adverse selection costs on
liquidity providers. Under a flat fee, the fee may be insufficient to
compensate for this adverse selection, causing liquidity providers to withdraw
and reducing market depth.

The dynamic fee addresses this through an automatic mechanism: when volatility
increases (indicating elevated informed trading activity), the fee increases
proportionally. This has two effects:

\begin{enumerate}
  \item \textbf{Direct compensation.} Higher fees directly offset the adverse
    selection cost borne by liquidity providers.
  \item \textbf{Deterrence.} The increased fee raises the cost of informed
    trading, reducing the profitability of frontrunning and sandwich attacks.
    This is particularly relevant in light of the MEV dynamics identified
    by~\citet{daian2020flash}.
\end{enumerate}

Formally, let $\pi_{\text{informed}}(\sigma, f)$ denote the expected profit
of an informed trader given volatility $\sigma$ and fee rate $f$. Under the
assumption that informed trading profit is proportional to volatility:
\begin{equation}
  \pi_{\text{informed}}(\sigma, f) = k \cdot \sigma - f \cdot V_{\text{trade}}
  \label{eq:informed-profit}
\end{equation}

For the dynamic fee, the breakeven condition
$\pi_{\text{informed}}(\sigma, f(\sigma)) = 0$ yields:
\begin{equation}
  k \cdot \sigma = \left(f_{\text{base}} + f_{\text{dyn}} \cdot \frac{\sigma}{\sigma_{\max}}\right) \cdot V_{\text{trade}}
  \label{eq:breakeven}
\end{equation}

The dynamic component $f_{\text{dyn}} \cdot \sigma / \sigma_{\max}$ scales
with the very quantity that makes informed trading profitable, providing a
natural hedge that a flat fee cannot replicate.

\subsection{On-Chain Volatility Oracle}
\label{subsec:vol-oracle}

The volatility $\sigma$ is computed entirely from on-chain data, eliminating
dependence on external price oracles:

\begin{equation}
  \sigma = \sqrt{\frac{1}{N-1} \sum_{i=1}^{N} \left(\ln\frac{p_i}{p_{i-1}} - \bar{r}\right)^2}
  \label{eq:realized-vol}
\end{equation}

where $p_i$ is the $i$-th trade price recorded on-chain, $N$ is the number of
trades in the trailing 24-hour window, and $\bar{r}$ is the mean log-return.
A minimum of $N \geq 10$ trades is required; if fewer trades exist, the fee
defaults to $f_{\text{base}}$.

This oracle-free design is a deliberate architectural choice: by computing
volatility from the protocol's own trade data rather than relying on external
feeds, the system eliminates oracle manipulation as an attack vector (see
Section~\ref{sec:security}).
